\section{Next actions}

The ISO/IEC Directives~\cite{jtc1:directives} state that the Chair, in
consultation with the relevant officers, decides which course of action to take:

\vspace{1em}
\begin{tcolorbox}[
  colback=blue!8!white,
  colframe=blue,
  title={\textbf{Section 2.6.4 from ISO/IEC JTC1 Directives}},
  breakable,
  width=\textwidth,
  boxrule=0.5pt
  ]
[\ldots] 

The Chair of the committee, in consultation with the
Secretary/Committee Manager and the Project Leader or working group Convenor,
and the Office of the CEO, shall take one of the following courses of action:

[\ldots]
\end{tcolorbox}

The actions available when the DIS is not approved are:

\vspace{1em}
\begin{tcolorbox}[
  colback=blue!8!white,
  colframe=blue,
  title={\textbf{Section 2.6.4 from ISO/IEC JTC1 Directives}},
  breakable,
  width=\textwidth,
  boxrule=0.5pt
  ]
[\ldots] 

c) when the approval criteria of 2.6.3 are not met:
\begin{enumerate}
\item to circulate the revised enquiry draft for eight weeks voting (see 2.6.1), or
\item to circulate the revised committee draft for comments, or
\item to circulate a revised draft as a draft Technical Specification, or
\item to cancel the project, subject to a decision by the committee.
\end{enumerate}

In case 3), the project shall start at the approval stage
\end{tcolorbox}

We briefly analyze each option below:

\subsection{Revised draft for enquiry}

The ISO/IEC directives specifically state \emph{``to circulate the revised
enquiry draft for eight weeks voting''} which means to revise the text of the
draft international standard after \emph{``at a meeting, address any comments
received and prepare, as needed, a revised version''}.

WG21 might (as suggested by~\cite{p4433r0}) decide not to make changes in
response to the requests from National Bodies to remove contracts from the
standard.
However, not doing anything at all and expecting that positions of National
Bodies will change does not seem a very likely scenario.

In our opinion, if WG21 decides not to make any change regarding contracts, the
ballot may produce a similar result.

There is another possible course of action, though. If WG21 removes the
contracts feature from the draft international standard, that would certainly
address objections directed specifically at contracts. This could reduce
opposition and improve the prospects for approval, although it would not
guarantee approval or increased consensus.

\textbf{Note}: The document \emph{SoV and Collated Comments}~\cite{sov:comments}
explicitly states in its text \emph{``The DIS ballot was not approved. The
results and comments have been sent to the Project Editor and WG 21 to review,
prepared [sic] the final DoC, and provide an updated draft text for further
processing.''}

\subsection{Revised draft for comments}

The second option is \emph{``to circulate the revised committee draft for
comments''}. This is a variation of the previous option: after the revised draft
is discussed at a meeting, it is circulated as a committee draft for comments
instead. This would require National Bodies to submit a new set of comments
before a new DIS ballot could be initiated.

This variation of the previous option would introduce an additional delay, and
we see no clear advantage that would increase consensus among National Bodies.

\subsection{Revised draft as a TS}

The third option is to abandon the plan to publish a new standard and instead
publish the new version as a Technical Specification.

This option might seem like a palatable compromise, but it would not resolve
the outstanding concerns. Those concerns would remain unresolved in a TS
rather than an IS.

Moreover, publishing a TS instead of an International Standard could delay
incorporation of many planned features into the standard, even if they were
available in the TS. These include reflection, erroneous behavior, linear
algebra, debugging support, executors and senders/receivers, SIMD support,
standard library hardening, and parallel range algorithms.

\subsection{Project cancellation}

We do not consider this a desirable outcome. However, if no other option is
available, this could happen. In that case, the baseline would be C++23, and a
new project would have to start from there.
